% Generated draft: requires mathematical and source review.
\documentclass[11pt]{article}
\usepackage[margin=25mm]{geometry}
\usepackage{fontspec}
\setmainfont{Latin Modern Roman}
\usepackage{amsmath,amssymb,mathtools,unicode-math}
\setmathfont{Latin Modern Math}
\usepackage{enumitem,needspace,xurl,hyperref}
\setlength{\parindent}{0pt}
\setlength{\parskip}{5pt}
\setlength{\emergencystretch}{3em}
\newcommand{\N}{\mathbb{N}}
\newcommand{\Z}{\mathbb{Z}}
\newcommand{\Q}{\mathbb{Q}}
\newcommand{\R}{\mathbb{R}}
\newcommand{\C}{\mathbb{C}}
\newcommand{\F}{\mathbb{F}}
\newcommand{\A}{\mathbb{A}}
\newcommand{\PP}{\mathbb{P}}
\newcommand{\E}{\mathbb{E}}
\newcommand{\eps}{\varepsilon}
\newcommand{\dd}{\,\mathrm{d}}
\begin{document}
\section*{UnsolvedMath ID: 30000903}
\textbf{Source problem code:} catalogue entry
\subsection{Definitions and mathematical statement}
Let $P$ be a convex $d$-polytope, not necessarily simplicial. The barycentric subdivision $\Delta=\operatorname{sd}(\partial P)$ is the simplicial complex whose vertices are nonempty proper faces of $P$ and whose simplices are chains of such faces. If $f_{i-1}$ counts its $(i-1)$-simplices, with $f_{-1}=1$, its $h$-polynomial is
\[
h_\Delta(t)=\sum_{i=0}^{d}f_{i-1}t^i(1-t)^{d-i}.
\]
The first conjecture asks whether every complex of this form has $h_\Delta(t)$ with only real zeros.

The second question concerns an arbitrary polynomial
\[
h(t)=\sum_{i=0}^{d}h_i t^i\in\Z[t],\qquad h_0=1,\quad h_i\geq1,
\]
with all zeros real, and asks whether its coefficient sequence satisfies the classical $g$-theorem conditions. In their usual full form these require symmetry $h_i=h_{d-i}$ and that
\[
g=(1,h_1-h_0,\ldots,h_{\lfloor d/2\rfloor}-h_{\lfloor d/2\rfloor-1})
\]
be an $M$-sequence, equivalently a Hilbert function of a standard graded Artinian algebra. In Macaulay form, its entries are nonnegative and obey $g_{i+1}\leq g_i^{\langle i\rangle}$: if $a=\binom{a_i}{i}+\cdots+\binom{a_j}{j}$ is the usual strictly decreasing binomial expansion, then $a^{\langle i\rangle}=\binom{a_i+1}{i+1}+\cdots+\binom{a_j+1}{j+1}$.

The general polynomial hypotheses do not include symmetry, and do not force it: $1+3t+2t^2=(1+t)(1+2t)$ satisfies them but has $h_2\ne h_0$. Hence a literal implication to the full classical $g$-theorem conditions needs a missing symmetry hypothesis or a weaker intended meaning of those conditions. The original record mentions unimodality and Macaulay first-difference inequalities without stating symmetry. Both questions and that discrepancy are retained; symmetry is not silently added, nor is the second question declared equivalent to the first.
\subsection{Short English statement}
Ask for real-rootedness of barycentric boundary $h$-polynomials of all convex polytopes. Separately, ask which $g$-theorem inequalities follow from positive integral real-rooted coefficients; the full classical conditions also require symmetry, absent from the supplied general question.
\subsection{Sources}
\begin{itemize}
\item Supplied report: \emph{Combinatorics} (2008), Oberwolfach Reports, DOI 10.4171/owr/2008/01, pp. 5--78; proposed year 2008. \url{https://doi.org/10.4171/owr/2008/01}.
\item Supplied Brenti--Welker reference on barycentric subdivisions: \url{https://arxiv.org/abs/math/0606356}.
\item Supplied Athanasiadis reference on local $h$-polynomials: \url{https://arxiv.org/abs/2402.06219}.
\item Local UnsolvedMath registry, permanent ID 30000903, catalogue entry; created 25 August 2026, updated 16 September 2026. Its corrected\_verified note reports comparison with the official report. The supplied 21 August 2026 triage labels the combined record partially solved, reporting simplicial/Cohen--Macaulay and selected shellable cases while retaining the general nonsimplicial-polytope question. The original separately describes the second conjecture using unimodality and Macaulay inequalities. This dated assessment is not independently updated here.
\end{itemize}
\end{document}
